\documentclass{article}
\usepackage{amsmath,amssymb,amsthm}
\usepackage{algorithm}
\usepackage[noend]{algpseudocode}
\usepackage{enumitem}
\usepackage{hyperref}
\usepackage{geometry}
\geometry{margin=1in}
\newtheorem{theorem}{Theorem}
\newtheorem{lemma}{Lemma}
\newcommand{\E}{\mathbb{E}}
\newcommand{\R}{\mathbb{R}}
\newcommand{\norm}[1]{\left\lVert#1\right\rVert}
\newcommand{\ip}[2]{\langle #1,#2\rangle}
\newcommand{\diag}{\mathrm{diag}}
\begin{document}

\section*{Randomized Coordinate Descent (RCD) for Non‑Convex Smooth Functions}

\section*{Mathematical Formulation}
Let $f:\R^{d}\rightarrow\R$ be differentiable and possibly non‑convex.  
For each coordinate $i\in\{1,\dots ,d\}$ we are given a smoothness constant $L_i>0$ such that

\[
\forall w\in\R^{d},\;\forall \alpha\in\R,\qquad 
f(w+\alpha e_i)\le f(w)+\alpha \partial_i f(w)+\frac{L_i}{2}\alpha^{2},
\tag{1}
\]
where $e_i$ is the $i$‑th canonical basis vector and $\partial_i f(w)=\frac{\partial f}{\partial w_i}(w)$.

At iteration $t$ a coordinate $i_t$ is drawn independently from a fixed distribution
\[
\Pr(i_t=i)=p_i,\qquad p_i>0,\;\sum_{i=1}^{d}p_i=1 .
\]

Define the \emph{scaled stochastic gradient}
\[
\tilde g_t\;:=\; \frac{1}{p_{i_t}}\,\partial_{i_t}f(w_t)\,e_{i_t}\in\R^{d}.
\tag{2}
\]
Then $\E[\tilde g_t\mid w_t]=\nabla f(w_t)$ (unbiasedness).

With a stepsize (learning rate) $\eta>0$ the update is
\[
w_{t+1}= w_t-\eta \tilde g_t .
\tag{3}
\]
Equivalently, only the selected coordinate is modified:
\[
w_{t+1}^{(i)}=
\begin{cases}
w_t^{(i)}-\displaystyle\frac{\eta}{p_i}\,\partial_i f(w_t), & i=i_t,\\[1ex]
w_t^{(i)}, & i\neq i_t .
\end{cases}
\tag{4}
\]

\section*{Pseudocode}
\begin{algorithm}[H]
\caption{Randomized Coordinate Descent (RCD)}\label{alg:RCD}
\begin{algorithmic}[1]
\Require Initial point $w_0\in\R^{d}$, stepsize $\eta>0$, probability vector $p\in\Delta^{d}$.
\For{$t=0,1,\dots,T-1$}
    \State Sample $i_t\sim\text{Multinomial}(p)$
    \State Compute partial derivative $g^{(i_t)}_t:=\partial_{i_t}f(w_t)$
    \State Form scaled gradient $\tilde g_t:=\frac{1}{p_{i_t}}g^{(i_t)}_t\,e_{i_t}$
    \State Update $w_{t+1}=w_t-\eta \tilde g_t$
\EndFor
\State \textbf{return} $w_T$ (or optionally the iterate with smallest $ \|\nabla f(w_t)\|$ )
\end{algorithmic}
\end{algorithm}

\section*{Dry Run Example}
Consider the one‑dimensional function $f(w)=(w-3)^2$.
Here $d=1$, $p_1=1$, $L_1=2$, and the gradient is $\partial_1 f(w)=2(w-3)$.
Take $w_0=0$, stepsize $\eta=0.1$.

\begin{tabular}{c|c|c|c}
$t$ & $w_t$ & $g_t^{(1)}=\partial_1 f(w_t)$ & $f(w_t)$\\\hline
0 & $0$ & $2(0-3)=-6$ & $9$\\
1 & $w_1 = 0-\frac{0.1}{1}(-6)=0+0.6=0.6$ & $2(0.6-3)=-4.8$ & $(0.6-3)^2=5.76$\\
2 & $w_2 = 0.6-\frac{0.1}{1}(-4.8)=0.6+0.48=1.08$ & $2(1.08-3)=-3.84$ & $(1.08-3)^2=3.6864$\\
3 & $w_3 = 1.08-\frac{0.1}{1}(-3.84)=1.08+0.384=1.464$ & $2(1.464-3)=-3.072$ & $(1.464-3)^2=2.3660$
\end{tabular}

The objective value decreases monotonically, illustrating the algorithm.

\newpage
\section*{Assumptions}
\begin{enumerate}[label=\textbf{A\arabic*.}, leftmargin=*]
\item \textbf{Coordinate-wise $L_i$‑smoothness.}  Inequality \eqref{1} holds for each $i$ with constant $L_i>0$.
\item \textbf{Bounded variance of the scaled estimator.} There exists $\sigma^2\ge0$ such that
\[
\E\!\left[\norm{\tilde g_t-\nabla f(w_t)}^{2}\mid w_t\right]\le\sigma^{2}.
\tag{5}
\]
\item \textbf{Finite lower bound.} $f^{\inf}:=\inf_{w\in\R^{d}}f(w)>-\infty$.
\item \textbf{Stepsize.} $\eta$ satisfies $0<\eta\le \min_{i}\frac{p_i}{L_i}$.
\end{enumerate}

\section*{Main Convergence Theorem}
\begin{theorem}[Expected Gradient Norm Decay]\label{thm:main}
Under Assumptions A1–A4, for any integer $T\ge1$ the iterates generated by Algorithm~\ref{alg:RCD} satisfy
\[
\frac{1}{T}\sum_{t=0}^{T-1}\E\!\left[\norm{\nabla f(w_t)}^{2}\right]
\;\le\;
\frac{2\bigl(f(w_0)-f^{\inf}\bigr)}{\eta\,T}
\;+\;
2\eta\,\sigma^{2}.
\tag{6}
\]
Consequently, choosing $\eta = \Theta\!\bigl(1/\sqrt{T}\bigr)$ yields
\[
\min_{0\le t<T}\E\!\left[\norm{\nabla f(w_t)}^{2}\right]=\mathcal{O}\!\left(\frac{1}{\sqrt{T}}\right).
\]
\end{theorem}

\section*{Theoretical Properties}
\begin{itemize}
\item \textbf{Stability.} The stepsize condition $0<\eta\le\min_i p_i/L_i$ guarantees that the expected one‑step decrease in $f$ is non‑negative (see Lemma~\ref{lem:descent}).
\item \textbf{Probability weighting.} The factor $1/p_i$ in \eqref{2} makes $\tilde g_t$ an unbiased estimator of the full gradient, which is crucial for obtaining \eqref{6}. If the weighting were omitted, the bound would involve the weighted average $\sum_i p_i L_i$, degrading the rate.
\item \textbf{Comparison to full‑gradient SGD.} The rate \eqref{6} matches the classical $\mathcal{O}(1/\sqrt{T})$ bound for non‑convex SGD, up to the constant $p_{\min}^{-1}$ hidden in $\eta$.
\item \textbf{Special case $p_i=1/d$ (uniform sampling).} The stepsize bound becomes $\eta\le \frac{1}{d\max_i L_i}$, and the convergence constant scales linearly with $d$.
\end{itemize}

\section*{Discussion}
The theorem shows that even without convexity, randomized coordinate descent converges to a point whose gradient norm vanishes in expectation at the same sub‑linear rate as stochastic gradient methods. The proof hinges on the coordinate‑wise smoothness and the unbiasedness created by the $1/p_i$ scaling. The result can be extended to block‑coordinate updates by grouping coordinates and using block‑wise smoothness constants.

\newpage
\section*{Preliminaries (Definitions and Lemmas)}
\begin{lemma}[One‑step Expected Descent]\label{lem:descent}
For any $w\in\R^{d}$ and any admissible stepsize $\eta$,
\[
\E\!\bigl[f(w_{+})\mid w\bigr]
\le f(w)-\frac{\eta}{2}\norm{\nabla f(w)}^{2}
+\frac{\eta^{2}}{2}\,\sigma^{2},
\tag{7}
\]
where $w_{+}=w-\eta\tilde g$ with $\tilde g$ defined in \eqref{2}.
\end{lemma}
\begin{proof}
Fix $w$ and denote $i\sim p$. Using coordinate‑wise smoothness \eqref{1} with $\alpha=-\eta\tilde g^{(i)}$,
\[
\begin{aligned}
f(w-\eta\tilde g)
&= f\!\bigl(w-\tfrac{\eta}{p_i}\partial_i f(w)e_i\bigr)\\
&\le f(w)-\frac{\eta}{p_i}\partial_i f(w)^{2}
+\frac{L_i}{2}\Bigl(\frac{\eta}{p_i}\partial_i f(w)\Bigr)^{2}.
\end{aligned}
\tag{8}
\]
Take expectation w.r.t.\ $i$:
\[
\begin{aligned}
\E[f(w_{+})\mid w]
&\le f(w)-\eta\sum_{i=1}^{d}p_i\frac{1}{p_i}\partial_i f(w)^{2}
+\frac{\eta^{2}}{2}\sum_{i=1}^{d}p_i\frac{L_i}{p_i^{2}}\partial_i f(w)^{2}\\
&= f(w)-\eta\norm{\nabla f(w)}^{2}
+\frac{\eta^{2}}{2}\sum_{i=1}^{d}\frac{L_i}{p_i}\partial_i f(w)^{2}.
\end{aligned}
\tag{9}
\]
Because of the stepsize restriction $\eta\le \min_i p_i/L_i$, we have $\frac{L_i}{p_i}\le\frac{1}{\eta}$, thus
\[
\frac{\eta^{2}}{2}\sum_{i}\frac{L_i}{p_i}\partial_i f(w)^{2}
\le\frac{\eta}{2}\norm{\nabla f(w)}^{2}.
\tag{10}
\]
Insert \eqref{10} into \eqref{9}:
\[
\E[f(w_{+})\mid w]\le f(w)-\frac{\eta}{2}\norm{\nabla f(w)}^{2}.
\tag{11}
\]
Finally, replace $w_{+}=w-\eta\tilde g$ by $w-\eta\tilde g+\eta(\tilde g-\nabla f(w))$ and use $\E[\tilde g-\nabla f(w)\mid w]=0$ together with variance bound \eqref{5} to obtain the extra $\frac{\eta^{2}}{2}\sigma^{2}$ term, yielding \eqref{7}.  ∎
\end{proof}

\section*{Main Proof (Step‑by‑Step Derivation)}
We prove Theorem~\ref{thm:main}.  
Let $w_t$ be the random iterate after $t$ updates and define $F_t:=\E[f(w_t)]$.

\begin{enumerate}[label=[Step \arabic*], leftmargin=*]
\item \textbf{Apply Lemma~\ref{lem:descent}.}  
Taking total expectation in \eqref{7} gives
\[
F_{t+1}\le F_t-\frac{\eta}{2}\E\!\left[\norm{\nabla f(w_t)}^{2}\right]+\frac{\eta^{2}}{2}\sigma^{2}.
\tag{12}
\]

\item \textbf{Telescoping sum.}  
Sum \eqref{12} from $t=0$ to $T-1$:
\[
F_T\le F_0-\frac{\eta}{2}\sum_{t=0}^{T-1}\E\!\left[\norm{\nabla f(w_t)}^{2}\right]
+\frac{\eta^{2}}{2}\sigma^{2}T.
\tag{13}
\]

\item \textbf{Lower bound $F_T$.}  
Since $f$ is bounded below by $f^{\inf}$,
\[
F_T\ge f^{\inf}.
\tag{14}
\]

\item \textbf{Rearrange to isolate the gradient sum.}  
Combine \eqref{13}–\eqref{14}:
\[
\frac{\eta}{2}\sum_{t=0}^{T-1}\E\!\left[\norm{\nabla f(w_t)}^{2}\right]
\le F_0-f^{\inf}+\frac{\eta^{2}}{2}\sigma^{2}T.
\tag{15}
\]

\item \textbf{Divide by $\eta T$.}  
\[
\frac{1}{T}\sum_{t=0}^{T-1}\E\!\left[\norm{\nabla f(w_t)}^{2}\right]
\le \frac{2\bigl(F_0-f^{\inf}\bigr)}{\eta T}+ \eta\sigma^{2}.
\tag{16}
\]

\item \textbf{Replace $F_0$ by $f(w_0)$.}  
Because $F_0=\E[f(w_0)]=f(w_0)$ (deterministic start), we obtain exactly the bound \eqref{6} stated in the theorem.
\end{enumerate}
Thus Theorem~\ref{thm:main} holds. ∎

\section*{Rate Derivation (With Constants)}
Choose $\eta = \sqrt{\dfrac{2\bigl(f(w_0)-f^{\inf}\bigr)}{\sigma^{2}T}}$, which respects the stepsize restriction for all sufficiently large $T$. Substituting into \eqref{6} yields
\[
\frac{1}{T}\sum_{t=0}^{T-1}\E\!\left[\norm{\nabla f(w_t)}^{2}\right]
\le 2\sqrt{\frac{2\bigl(f(w_0)-f^{\inf}\bigr)\sigma^{2}}{T}}
= \mathcal{O}\!\left(\frac{1}{\sqrt{T}}\right).
\tag{17}
\]

\section*{Special Cases}
\begin{enumerate}[label=\alph*)]
\item \textbf{Uniform sampling.} $p_i = 1/d$ gives $\eta\le \frac{1}{d\max_i L_i}$ and the bound \eqref{6} scales with $d$.
\item \textbf{Importance sampling.} Selecting $p_i \propto L_i$ maximizes the admissible stepsize $\eta\le \bigl(\sum_i L_i\big)^{-1}$, often leading to a tighter bound.
\item \textbf{Deterministic cyclic coordinate descent.} Setting $p_i=1$ for a single coordinate at a time (i.e., $d=1$ block) recovers the classic non‑convex coordinate descent bound.
\end{enumerate}

\end{document}